\section{总结与展望}

之间取得平衡（图 4-10）。




\begin{figure}[htbp]
  \centering
  \includegraphics[width=0.78\textwidth]{pic/fig4-10.png}
  \caption{多种方法对接平稳度对比}
\end{figure}


4.5.4 对比结果总结
综合上述分析，本文方法的主要特点并不是在所有指标上优于人工或其他算法，
而是在特定研究阶段和特定应用目标下具有较好的适用性。其优势主要体现在以下几

个方面：
第一，本文方法能够把机场牵引作业拆解为多个可控阶段，使“前往飞行
器区域、对接飞行器、推出飞行器、解除牵引、车辆撤离”形成较完整的仿真
流程。
第二，ROS 与 Gazebo 提供了可重复测试环境，便于反复检查车辆模型、控
制策略和作业流程；相比直接开展真实车辆测试，试错成本和安全风险更低。
第三，相对位姿误差被引入控制过程后，车辆能够根据目标距离、横向偏
差和航向偏差修正动作，比单纯开环脚本更能体现自动对接的反馈特征。
第四，感知、任务管理、局部目标生成、误差控制和执行层被分开组织，
后续若增加传感器融合、目标识别或路径优化算法，可以在相应节点上继续扩
展。
\begin{figure}[htbp]
  \centering
  \includegraphics[width=0.78\textwidth]{pic/fig4-11.png}
  \caption{不同方法的综合能力对比}
\end{figure}

总体而言，本文方法是在人工牵引、开环脚本和常见路径跟踪算法之间形
成的折中方案。它主要服务于 ROS-Gazebo 仿真验证，在有限实验条件下提高了






流程完整性、可重复性和任务鲁棒性。





4.6 本章小结
本章围绕 ROS-Gazebo 联合仿真系统展开，说明了软件栈、牵引车模型、传
感器挂载、任务模式管理以及接口关系。通过 URDF 模型、Gazebo 场景、
ros\_control 控制器和 ROS 节点通信，牵引车能够完成模型加载、转向回中、
轮速控制、目标接近和停车保持等动作，为第三章方法提供验证环境。后续可
替换车辆模型、调整控制器参数，或增加更复杂的机场地图和传感器配置。
5.总结与展望
\subsection{总结}
本文围绕机场飞机牵引与自动对接作业需求，以自动牵引车和飞机协同控
制为研究对象，基于 ROS 与 Gazebo 联合仿真平台，开展了作业问题抽象、车辆
运动学建模、感知与控制方法设计以及仿真实现等方面的研究。
在建模方面，本文把飞机推出和前起落架对接过程整理为低速、近距和多
约束的车辆控制问题，并建立双桥转向牵引车的平面运动学模型。通过比较前
桥单转、前后桥反向和同向转向等模式，明确速度、转角和曲率之间的约束关
系，为后续控制设计提供基础。



